\section{Little Things}

Now that we can create simple scale maps for any tune, we can take the next step in our improvisation journey. With our scale maps we are essentially using four kinds of scales. The major scale which is used for when we have a major chord with tonic function, the minor 7 scale for when we have a minor seventh chord with predominant function, the minor scale with the raised seventh for when we have a minor chord with tonic function, and the dominant scale which is used for chords with dominant function. Next, we will practice what we will call \emph{little things}. These are short melodic modules that we will practice in order to develop the technical dexterity required to improvise fluidly on tunes.

The first few things we learn will be familiar. If the scale maps exercises are about learning to play the scales up and down, then a large part of our study of little things will be about finding different ways to practice our scales. One simple thing to try is playing our scale in thinds up and down. 

\bookExample{Scales in Thirds}{./excerpts/in-thirds.ly}

 Practicing this in all 12 keys with the four different scales will take quite some time to get under our fingers. Even with this simple example, there a number of small variations we can practice.

\begin{bookStyledBulletList}{Variations on the Thirds}
	\item Higher note first.
	\item Offset the rhythm by one 8th note.
	\item Play slower or faster rhythmic values.
	\item Add a whole step above or below one of the notes.
	\item Add a half step above or below one of the notes.
	\item Change the rhythm
	\item Add some bebop turns. (16th note triplets ala Parker)
\end{bookStyledBulletList}

Once we are comfortable with thirds like the above in all the necessary keys, we can apply this to a tune. It's important to remember to play this in time with a metronome. If we are comfortable performing the scale map as normal, we can now play through our tune in thirds  without jumping around, like the scale maps do. We should maintain our direction and simply switch scales instead of jumping to the tonic of the other scale. This can become quite complex if we add some of the variations listed above to this process.
